\documentclass{chsmc}
\chsmcsetup{
	year = 2011,
	part = 1,
	type = A,
	show-answers = false,
}
\begin{document}

\section{填空题：本大题共 8 小题，每小题 8 分，共 64 分.}

\begin{enumerate}
	\item 设集合 $A = \{a_1,a_2,a_3,a_4\}$，若 $A$ 中所有三元子集的三个元素之和组成的集合为
		$B = \{-1,3,5,8\}$，则集合 $A =$ \blankline
	\item 函数 $f(x) = \dfrac{\sqrt{x^2+1}}{x-1}$ 的值域为 \blankline
	\item 设 $a$, $b$ 为正实数，$\dfrac{1}{a} + \dfrac{1}{b} \leqslant 2\sqrt{2}$，
		$(a-b)^2 = 4(ab)^3$，则 $\log_a b =$ \blankline
	\item 如果 $\cos^5\theta - \sin^5\theta < 7(\sin^3\theta - \cos^3\theta)$，
		$\theta \in [0,2\pi)$，那么 $\theta$ 的取值范围是 \blankline
	\item 现安排 $7$ 名同学去参加 $5$ 个运动项目，要求甲、乙两同学不能参加同一个项目，
		每个项目都有人参加，每人只参加一个项目，则满足上述要求的不同安排方案数为
		\blankline（用数字作答）
	\item 在四面体 $ABCD$ 中，已知 $\angle ADB = \angle BDC = \angle CDA = 60^\circ$，
		$AD = BD = 3$, $CD = 2$，则四面体 $ABCD$ 的外接球的半径为 \blankline
	\item 直线 $x - 2y - 1 = 0$ 与抛物线 $y^2 = 4x$ 交于 $A$, $B$ 两点，
		$C$ 为抛物线上的一点，$\angle ACB = 90^\circ$，则点 $C$ 的坐标为 \blankline
	\item 已知 $a_n = \mathrm{C}_{200}^n\cdot (\sqrt[3]{6})^{200-n}\cdot
		\left(\dfrac{1}{\sqrt{2}}\right)^n$（$n = 1,2,\cdots,95$），
		则数列 $\{a_n\}$ 中整数项的个数为 \blankline
\end{enumerate}

\section{解答题：本大题共 3 个小题，满分 56 分. 解答应写出文字说明、证明过程或演算步骤.}

\begin{enumerate}[resume]
	\item (本题满分 16 分) 设函数 $f(x) = |\lg(x+1)|$，实数 $a$, $b$（$a < b$）满足
		$f(a) = f\left(\dfrac{b+1}{b+2}\right)$，$f(10a + 6b + 21) = 4\lg 2$，求 $a$, $b$
		的值.
	\item (本题满分 20 分) 已知数列 $\{a_n\}$ 满足：$a_1 = 2t - 3$
		（$t \in \mathbf{R}$ 且 $t \neq \pm 1$），
		$a_{n+1} = \dfrac{(2t^{n+1}-3)a_n + 2(t-1)t^n - 1}{a_n + 2t^n - 1}$
		（$n \in \mathbf{N}^*$）.
    \begin{enumerate}[(1)]
      \item 求数列 $\{a_n\}$ 的通项公式；
      \item 若 $t > 0$，试比较 $a_{n+1}$ 与 $a_n$ 的大小.
    \end{enumerate}
	\item (本题满分 20 分) 作斜率为 $\dfrac{1}{3}$ 的直线 $l$ 与椭圆
		$C\colon \dfrac{x^2}{36} + \dfrac{y^2}{4} = 1$ 交于 $A$, $B$ 两点（如图所示），
		且 $P(3\sqrt{2},\sqrt{2})$ 在直线 $l$ 的左上方.
    \begin{enumerate}[(1)]
      \item 证明：$\triangle PAB$ 的内切圆的圆心在一条定直线上；
      \item 若 $\angle APB = 60^\circ$，求 $\triangle PAB$ 的内切圆的面积.
    \end{enumerate}
\end{enumerate}

\end{document}
